% =============================================================
% Appendix: Case Studies for Pattern Analysis
%
% Required preamble (add once, near the top of your main file):
%   \usepackage[most]{tcolorbox}
%
% Design notes for two-column layouts (NeurIPS/ICLR/etc.):
%   - All tcolorboxes use width=\linewidth and breakable so they
%     stay inside a single column and can split across columns.
%   - Per-model headings are folded into the box title to avoid
%     run-in \paragraph{...} interacting badly with tcolorbox
%     placement.
% =============================================================
\section{Limitations}
\label{sec:limitations}

Several limitations of the present work are acknowledged.
First, the four composition operators address representative patterns but do not fully capture the entire design space.
More complex structures, such as conditional branching and bounded loops, remain unexplored.
Second, the effectiveness of RACES relies on the underlying model demonstrating adequate reasoning competence. Models with limited capabilities typically yield sparse rewards on composite tasks. Finally, composite environments generate extended reasoning traces, which require a sufficiently large context window (32K in these experiments) to enable effective training.
These limitations present opportunities for future research rather than imposing fundamental constraints on the framework.


\section{Declaration of LLM Usage}
\label{sec:llm_declaration}

\paragraph{Methodological use.} Following the SCALER~\citep{xu2026scalersyntheticscalableadaptivelearning} and RESYN~\citep{he2026resyn} synthesis paradigms, Claude-Sonnet-4.5 generates candidate verifiable environments, which subsequently pass through our two-stage filtering (static code self-consistency followed by multi-sample output consistency) before entering the initial environment pool.
\paragraph{Writing assistance.} LLMs were used exclusively for limited writing assistance, specifically for grammar checking and minor sentence-level refinement.

\section{Broader Impacts}
\label{sec:broader_impacts}

\paragraph{Positive impacts.} RACES advances automatic verifiable RL data construction, reducing reliance on human annotation. By transforming a small pool of environments into a structurally diverse and effectively unbounded space of composite tasks, RACES increases the accessibility of verifiable RL for research groups lacking industrial-scale annotation resources.

\paragraph{Potential negative impacts.} As with most foundational work on improving LLM capabilities, the techniques that enhance reasoning may also be misapplied. Specifically, the methodology could be adapted to synthesize verifiable training data for malicious purposes, and the energy consumption associated with large-scale RL training may contribute to environmental impact. Since model checkpoints are not released, the most direct vector of misuse is absent in this work. The remaining artifacts serve as general-purpose research infrastructure rather than capability-amplifying weights. The 300 environments do not contain content related to poison or dangerous material.

\section{Initial environments construction.}
\label{init_envs}
RACES constructs an initial pool $\mathcal{E}$ of 300 environments, sourced from three categories: (1) standardized algorithmic problem datasets~\citep{xia2025leetcodedatasettemporaldatasetrobust}\footnote{45 of the 300 environments are derived from algorithms.}; (2) environments automatically generated by Claude-Sonnet-4.5, following the synthesis paradigms of SCALER~\citep{xu2026scalersyntheticscalableadaptivelearning} and RESYN~\citep{he2026resyn}; and (3) manually authored environments.
Each environment undergoes a quality assurance review. For each, 20 inputs are sampled from $G_e$, and Claude-Sonnet-4.5 is tasked with solving $D_e(x)$ multiple times, requiring a $V_e$-judged pass rate exceeding 95%. For environments with a low pass rate, $D_e$ is manually refined.
\label{app:quality}

This section outlines the filtering criteria applied to each candidate extension during composition-path discovery.

\paragraph{Executability.} Each candidate is executed on the current state within a sandboxed process and is discarded if it raises an exception, exceeds a 2-second wall-clock timeout, or returns no output. Extensions requiring more than 400 atomic steps are also rejected, as this is considered beyond the reasonable effort of a human solver.

\paragraph{Non-degeneracy.} Let $y_{t+1}$ denote the resulting output and $s = \mathrm{str}(y_{t+1})$ its string representation. An extension is discarded if $y_{t+1}$ (i) repeats an earlier state in the same path or duplicates a sibling output at the same frontier, (ii) equals $0$ or $1$, or is an integer greater than $500$, (iii) has a length exceeding $100$, or (iv) consists of a single repeated character.

After path discovery, paths with a length less than two are discarded, as single-step paths reduce to environments.

\section{Case Studies for Pattern Analysis}
\label{app:cases}

This appendix presents the prompts and key response excerpts for the three cases referenced in Section~\ref{sec:pattern_analysis}.
The excerpts focus on the decision points where the three models diverge; complete responses are released in the supplementary materials. We label the models as Base, $RL_{individual}$, and $RL_{RACES}$ for consistency with the main text.

% -------------------------------------------------------------
\subsection{AIME: \texorpdfstring{$2\times2$}{2x2} Grid-Coloring}
\label{app:case-grid-coloring}

\textbf{Pattern illustrated.} A shift from ad-hoc enumeration toward representation-mediated counting: $RL_{RACES}$ introduces a reusable function abstraction that decouples argument tracking from arithmetic and verifies the final sum in two independent ways.

\smallskip
\begin{tcolorbox}[
  enhanced, breakable, width=\linewidth,
  colback=gray!5, colframe=gray!50,
  title={Problem (AIME 2025), ground truth: $\boxed{82}$},
  fonttitle=\bfseries\sffamily, coltitle=black
]
Four unit squares form a $2 \times 2$ grid. Each of the 12 unit line
segments forming the sides of the squares is colored either red or
blue in such a way that each unit square has 2 red sides and 2 blue
sides. Find the number of such colorings.
\end{tcolorbox}

\smallskip
\begin{tcolorbox}[
  enhanced, breakable, width=\linewidth,
  colback=red!3, colframe=red!40,
  title={\textbf{Base}, variable-confusion in case-table filling},
  fonttitle=\sffamily, coltitle=black
]
\begin{Verbatim}[fontsize=\footnotesize, breaklines, breakanywhere]
[Sets up grid, labels edges, identifies 4 internal + 8 boundary edges correctly.]

Each square must have exactly 2 red edges. [...] For each square, given the internal edges, the red edges in the square are determined by the pair. [...]

We can consider the possible red edge configurations for each square and see when the shared edges match. [...]

[Derives four equations:]
- (1): $x_1 + x_5 = 2 - a - c$, i.e. $s_1 = 2 - a - c$
- (2): $x_2 + x_7 = 2 - b - c$
- (3): $x_3 + x_6 = 2 - a - d$
- (4): $x_4 + x_8 = 2 - b - d$

[Builds a 16-row table of $(a,b,c,d)$ values:]

| a | b | c | d | s1 | s2 | s3 | s4 | w1 | w2 | w3 | w4 | f |
|---|---|---|---|----|----|----|----|----|----|----|----|---|
|0|0|0|0|2|2|2|2|1|1|1|1|1|
|0|0|0|1|2|2|1|1|1|1|2|2|4|
|**0**|**0**|**1**|**0**|**1**|**1**|**1**|**1**|**2**|**2**|**2**|**2**|**16**|
[...]

**<- KEY ERROR:** For $(a,b,c,d)=(0,0,1,0)$, RAW writes $s_3=2-a-d=2-0-0=1$ and $s_4=2-b-d=2-0-0=1$. The correct values are $s_3=2$ and $s_4=2$ (since $d=0$), giving $f=4$ not $16$. This variable-confusion error recurs in $\geq$8 of 16 rows.

[Sums all $f$ values to get 158.]

$$\boxed{158}$$
\end{Verbatim}
\end{tcolorbox}

\smallskip
\begin{tcolorbox}[
  enhanced, breakable, width=\linewidth,
  colback=orange!3, colframe=orange!50,
  title={\textbf{$RL_{individual}$}, single inconsistent state update},
  fonttitle=\sffamily, coltitle=black
]
\begin{Verbatim}[fontsize=\footnotesize, breaklines, breakanywhere]
[Uses case analysis by number of red internal edges $i_r \in \{0,1,2,3,4\}$.]

**Case $i_r=0$:** All internal blue. Each square forces both boundary edges red. -> **1 coloring** $\checkmark$

**Case $i_r=1$:** [Correctly derives constraints for H3 red:] $h1+v1=1$, $h2+v5=2$, $h5+v2=1$, $h6+v6=2$ -> $2 \times 2 = 4$ per choice, $\times 4$ choices -> **16** $\checkmark$

**Case $i_r=2$:**
- Opposite pairs: each gives $2^4=16$ -> $2 \times 16 = 32$ $\checkmark$
- Adjacent pairs: each gives $4$ -> $4 \times 4 = 16$ $\checkmark$
- Subtotal: **48** $\checkmark$

**Case $i_r=3$:** 4 choices $\times$ 4 each -> **16** $\checkmark$

**Case $i_r=4$:** All internal red.

[...]

- D: edges: H4 (red), H6, V4 (red), V6 -> red: $1 + h6 + **0** + v6$

**<- KEY ERROR:** Writes V4's contribution as $+0$ (treating it as blue), despite explicitly stating "V4 all red" at the start of this case. Correct: $1+h6+1+v6=2 \Rightarrow h6+v6=0 \Rightarrow$ **1 solution**, not $h6+v6=1 \Rightarrow$ 2 solutions.

[Sums: $1+16+48+16+2 = 83$]

$$\boxed{83}$$
\end{Verbatim}

\end{tcolorbox}

\smallskip
\begin{tcolorbox}[
  enhanced, breakable, width=\linewidth,
  colback=green!3, colframe=green!75!black,
  title={\textbf{RACES}, functional abstraction with two-way verification},
  fonttitle=\sffamily, coltitle=black
]
\begin{Verbatim}[fontsize=\footnotesize, breaklines, breakanywhere]

[Sets up same grid and equations, then introduces a key abstraction:]

For any two bits $x,y \in \{0,1\}$, define:
$$f(x,y) = \begin{cases} 1 & \text{if } x+y \in \{0,2\} \\ 2 & \text{if } x+y=1 \end{cases}$$

Then:
- $N_1 = f(a,c)$, $N_2 = f(b,c)$, $N_3 = f(a,d)$, $N_4 = f(b,d)$

Total = $\sum_{a,b,c,d} f(a,c) \cdot f(b,c) \cdot f(a,d) \cdot f(b,d)$

[Evaluates all 16 rows cleanly — each is a direct lookup of $f$:]

| a | b | c | d | f(a,c) | f(b,c) | f(a,d) | f(b,d) | Product |
|---|---|---|---|--------|--------|--------|--------|---------|
| 0 | 0 | 0 | 0 | 1 | 1 | 1 | 1 | 1 |
| 0 | 0 | 0 | 1 | 1 | 1 | 2 | 2 | 4 |
| 0 | 0 | 1 | 0 | 2 | 2 | 1 | 1 | **4** |
[... all 16 rows computed correctly ...]

[**Verification by grouping:**]
- 12 rows with product 4 -> $12 \times 4 = 48$
- 2 rows with product 16 -> $2 \times 16 = 32$
- 2 rows with product 1 -> $2 \times 1 = 2$
- Total: $48 + 32 + 2 = 82$ $\checkmark$

[**Sequential verification:**]
$1+4+4+16+4+4+4+4+4+4+4+4+16+4+4+1 = 82$ $\checkmark$

$$\boxed{82}$$

\end{Verbatim}
\end{tcolorbox}

% -------------------------------------------------------------
\subsection{Enigmata: List-Transformation Rule Induction}
\label{app:case-list-transform}

\textbf{Pattern illustrated.} Hypothesis verification against all
demonstrations, combined with explicit index--value tracking.
$RL_{RACES}$ systematically rules out candidate rules that fail on
earlier examples, where Base and $RL_{individual}$ repeatedly
propose rules contradicted by their own observations.

\smallskip
\begin{tcolorbox}[
  enhanced, breakable, width=\linewidth,
  colback=gray!5, colframe=gray!50,
  title={Problem, ground truth: $[66, 66, 66]$},
  fonttitle=\bfseries\sffamily, coltitle=black
]
Apply a function to the final input list to generate the output
list. Use any preceding inputs and outputs as examples.

\smallskip
\textbf{Examples:}\\
$[4, 34, 17, 7, 5] \to [17, 7, 5]$\\
$[2, 42, 1, 7, 4, 4, 59, 3, 28, 72] \to [7, 4, 4, 3, 28]$\\
$[88, 8, 55, 12, 4, 2, 75] \to [55, 12, 75]$\\
$[18, 40, 35, 57, 18, 18, 33, 35, 35] \to [35, 33, 35, 35]$

\smallskip
\textbf{Test:} $[46, 94, 66, 98, 66, 66] \to ?$
\end{tcolorbox}

\smallskip
\begin{tcolorbox}[
  enhanced, breakable, width=\linewidth,
  colback=red!3, colframe=red!40,
  title={\textbf{Base}, broad hypothesis exploration without verification},
  fonttitle=\sffamily, coltitle=black
]
\begin{Verbatim}[fontsize=\footnotesize, breaklines, breakanywhere]

[Tries dozens of hypotheses: local maxima removal, suffix extraction, runs of duplicates, threshold filtering...]

Perhaps the function is to **remove all elements that are strictly greater than the previous element?**

Test example 1: [4,34,17,7,5]
- 4->34: 34>4 -> remove. 17<34->keep. 7<17->keep. 5<7->keep.
So output: [17,7,5] — matches!

Check example 2: [...]
Kept: 2, 1, 4, 4, 3 — **not [7,4,4,3,28]** -> fails.

[Tries ~15 more hypotheses, none verified against all examples.]

[...] I think I have to give up and guess [...] based on the only run of duplicates at the end.

```
[66, 66]
```
\end{Verbatim}
\end{tcolorbox}

\smallskip
\begin{tcolorbox}[
  enhanced, breakable, width=\linewidth,
  colback=orange!3, colframe=orange!50,
  title={\textbf{$RL_{individual}$}, contradictions noted but not used},
  fonttitle=\sffamily, coltitle=black
]
\begin{Verbatim}[fontsize=\footnotesize, breaklines, breakanywhere]

[Follows nearly identical exploration. Tests same hypotheses.]

[...] In example 1: no duplicates -> **output should be empty? But output is [17,7,5]** — so not that.

**<- Notes the contradiction explicitly, but continues with the falsified rule anyway.**

[...] I think the function is to keep the last run of identical values of length at least 2. [...] For the test problem: only the last two 66s.

```
[66, 66]
```
\end{Verbatim}
\end{tcolorbox}

\smallskip
\begin{tcolorbox}[
  enhanced, breakable, width=\linewidth,
  colback=green!5, colframe=green!75!black,
  title={\textbf{RACES}, index tracking and rule preservation across all demos},
  fonttitle=\sffamily, coltitle=black
]
\begin{Verbatim}[fontsize=\footnotesize, breaklines, breakanywhere]
[Tests hypotheses systematically, verifying each against ALL examples before proceeding.]

[...] Example 1: from index 2 -> [17,7,5]. Example 2: last 5 elements... **not matching**. [Rejects suffix hypothesis.]

[Shifts to tracking indices explicitly:]

Example 4: output [35,33,35,35] -> indices 2,6,7,8
- 35 appears at indices 2, 7, 8 — appears 3 times total
- 33 appears at index 6 — once

[...] outputs consistently included values appearing multiple times in the input (66 appears at indices 2, 4, 5) -> all three occurrences of the repeated value.

```
[66, 66, 66]
```
\end{Verbatim}
\end{tcolorbox}

% -------------------------------------------------------------
\subsection{Enigmata: Sum Skyscraper Logic Puzzle}
\label{app:case-skyscraper}

\textbf{Pattern illustrated.} Stateful constraint propagation with
deliberate decomposition. $RL_{RACES}$ switches the search axis
(rows $\to$ columns) so that the smaller search space is exhausted
first, then propagates the resulting fixed values to constrain the
remaining cells; both Base and $RL_{individual}$ commit to an
incomplete row-wise search and incorrectly conclude that no solution
exists.

\smallskip
\begin{tcolorbox}[
  enhanced, breakable, width=\linewidth,
  colback=gray!5, colframe=gray!50,
  title={Problem, Sum Skyscraper $4{\times}4$},
  fonttitle=\bfseries\sffamily, coltitle=black
]
Fill a $4\times 4$ grid with heights $1$--$4$ so that each row and
column is a permutation of $\{1,2,3,4\}$ and the visibility-sum
clues on all four sides are satisfied (a clue gives the total height
of buildings visible from that side; taller buildings block shorter
ones behind them).

\smallskip
\textbf{Clues} (top, left, right, bottom):
\begin{verbatim}
7 9 4 5
7 9 4 5
5 4 7 9
5 4 7 9
\end{verbatim}

\smallskip
\textbf{Ground-truth answer:}
$\bigl[\,[3,2,4,1],\,[2,3,1,4],\,[4,1,2,3],\,[1,4,3,2]\,\bigr]$.
\end{tcolorbox}

\smallskip
\begin{tcolorbox}[
  enhanced, breakable, width=\linewidth,
  colback=red!3, colframe=red!40,
  title={\textbf{Base}, visibility assumption error in row-wise casework},
  fonttitle=\sffamily, coltitle=black
]
\begin{Verbatim}[fontsize=\footnotesize, breaklines, breakanywhere]
[Correctly identifies visibility-sum semantics. Attempts row-by-row search for Row 1 (left=7, right=5).]

[...tries ~15 permutations...]

Try [3,2,4,1]:
- Left: 3->visible, 2<3->blocked, 4>3->visible -> sum=3+4=7 $\checkmark$
- Right: 1->visible, 4>1->visible -> sum=1+4=5 $\checkmark$

**<- Actually finds a valid Row 1!** But then analyzes right-sum=5 in general:

Case 2: d=2, c=3: row [a,b,3,2], need left=7.
Available {a,b}={1,4}.

Option 1: a=1, b=4 -> left: 1+4=5  x 
Option 2: a=4, b=1 -> left: 4 only  x 

<- KEY ERROR: Ignores [3,2,4,1] found moments earlier. Concludes "no valid solution."
\end{Verbatim}
\end{tcolorbox}

\smallskip
\begin{tcolorbox}[
  enhanced, breakable, width=\linewidth,
  colback=orange!3, colframe=orange!50,
  title={\textbf{$RL_{individual}$}, incomplete row enumeration},
  fonttitle=\sffamily, coltitle=black
]

\begin{Verbatim}[fontsize=\footnotesize, breaklines, breakanywhere]

[Also searches row-by-row. Tests many permutations for Row 1 but never tests [3,2,4,1].]

[...] Try [3,1,4,2]: Left=7 $\checkmark$, Right=6  x 
Try [3,4,1,2]: Left=7 $\checkmark$, Right=6  x 

[Moves to Row 4 (right=9):] "No sequence gives sum = 9."

<- Fails to test [4,3,1,2]: right-to-left visibility: 2->3>2->4>3 -> sum=2+3+4=9. Concludes "No valid solution."

\end{Verbatim}
\end{tcolorbox}

\smallskip
\begin{tcolorbox}[
  enhanced, breakable, width=\linewidth,
  colback=green!3, colframe=green!75!black,
  title={\textbf{RACES}, column-wise decomposition with cascading propagation},
  fonttitle=\sffamily, coltitle=black
]
\begin{Verbatim}[fontsize=\footnotesize, breaklines, breakanywhere]

[**Key strategic choice:** decomposes column-by-column instead of row-by-row.]

**Column 1 (top=7, bottom=5):** Enumerates all 24 permutations systematically.

[...] (3,2,4,1): Top: 3+4=7 $\checkmark$. Bottom: 1->visible, 4>1->visible -> sum=1+4=5 $\checkmark$

**-> Unique valid assignment: A=3, E=2, I=4, M=1**

**Column 2 (top=9, bottom=4):** Bottom=4 forces N=4 (only tallest visible). Then:
- (2,1,3,4): top sum=2+3+4=9 $\checkmark$
- (2,3,1,4): top sum=2+3+4=9 $\checkmark$

**Column 3 (top=4):** Forces C=4 (only way to get sum=4). Then bottom=7 yields two options for (G,K,O).

**Column 4 (bottom=9):** Finds (1,4,3,2) and (4,3,1,2) both valid.

[**Constraint propagation:**]
Row 1: A=3, B=2, C=4 -> D must be **1** (only value left).
-> Fixes Column 4 to Option X: D=1, H=4, L=3, P=2.

Row 2: E=2, H=4 -> F,G $\in$ {1,3}. But E=2 already in row -> **G$\neq$2** -> G=1 -> F=3.
-> Fixes Column 3 Option A and Column 2 Option 2.

[**Full verification of all 16 clues** — all pass.]

```
3 2 4 1
2 3 1 4
4 1 2 3
1 4 3 2
```
\end{Verbatim}
\end{tcolorbox}